\section{Proposed Methodology}
\label{sec:methodology}

% \subsection{Design Motivation: Addressing Industrial Challenges}
\subsection{설계 동기: 산업적 과제 해결}
% The design of our system is motivated by three critical challenges observed in industrial inspection.
우리 시스템의 설계는 산업 점검에서 관찰된 세 가지 중요한 과제에 의해 동기 부여되었습니다.
% First, the \textbf{Geometric Distortion Challenge} arises from the fact that inspection robots cannot always align themselves perfectly with target equipment, necessitating robust rectification modules.
첫째, \textbf{기하학적 왜곡 과제}는 점검 로봇이 대상 장비와 항상 완벽하게 정렬될 수 없다는 사실에서 발생하며, 이에 따라 강건한 보정 모듈이 필요합니다.
% Second, the \textbf{Semantic Diversity Challenge} refers to the wide variety of inspection points; while gauges can be modeled geometrically, items like safety signs require high-level reasoning.
둘째, \textbf{의미론적 다양성 과제}는 점검 지점의 광범위한 다양성을 의미합니다. 게이지는 기하학적으로 모델링될 수 있지만, 안전 표지판과 같은 항목은 고수준의 추론이 필요합니다.
% Third, the \textbf{Efficiency-Accuracy Trade-off} requires a system that provides deep analysis without stalling the robot's movement.
셋째, \textbf{효율성-정확도 트레이드오프}는 로봇의 움직임을 멈추지 않으면서도 심층적인 분석을 제공하는 시스템을 요구합니다.
% Our modular methodology is specifically tailored to address these needs through a layered approach.
우리의 모듈형 방법론은 계층적 접근 방식을 통해 이러한 요구 사항을 해결하도록 특별히 설계되었습니다.

% The proposed industrial inspection system is designed as a multi-layered modular framework that bridges the gap between raw visual data and structured diagnostic reporting.
제안된 산업 점검 시스템은 원시 시각 데이터와 구조화된 진단 보고 사이의 간극을 메우는 다계층 모듈형 프레임워크로 설계되었습니다.
% The methodology is divided into four functional layers: the Orchestration Layer, the Perception and Pre-processing Layer, the Specialized Diagnostics Layer, and the Resource Management Layer.
이 방법론은 오케스트레이션 계층, 인지 및 전처리 계층, 특화 진단 계층, 그리고 자원 관리 계층의 네 가지 기능 계층으로 나뉩니다.
% The entire system flow is illustrated in Figure \ref{fig:architecture}.
전체 시스템 흐름은 그림 \ref{fig:architecture}에 예시되어 있습니다.

\begin{figure*}[t]
\centering
\fbox{\rule{0pt}{3in} \rule{0.95\linewidth}{0pt}} % Placeholder for detailed high-level architecture
% \caption{The high-level architecture of the integrated industrial inspection system. The system is divided into four main layers: (1) Orchestration Layer for task management, (2) Perception and Pre-processing Layer for object detection and geometric refinement, (3) Specialized Diagnostics Layer for domain-specific interpretation (AG, DG, LED, VLM), and (4) Resource Management Layer for performance optimization.}
\caption{통합 산업 점검 시스템의 고수준 아키텍처. 시스템은 네 개의 주요 계층으로 나뉩니다: (1) 작업 관리를 위한 오케스트레이션 계층, (2) 객체 탐지 및 기하학적 정제를 위한 인지 및 전처리 계층, (3) 도메인 특화 해석을 위한 특화 진단 계층 (AG, DG, LED, VLM), (4) 성능 최적화를 위한 자원 관리 계층.}
\label{fig:architecture}
\end{figure*}

% \subsection{Orchestration Layer: Task Lifecycle Management}
\subsection{오케스트레이션 계층: 작업 생명주기 관리}
% At the highest level, the system is governed by a central orchestration engine (\texttt{main.py} and the \texttt{DiagnosisSystem} class).
최상위 수준에서 시스템은 중앙 오케스트레이션 엔진(\texttt{main.py} 및 \texttt{DiagnosisSystem} 클래스)에 의해 제어됩니다.
% This layer is responsible for the end-to-end lifecycle of an inspection mission.
이 계층은 점검 미션의 엔드 투 엔드 생명주기를 책임집니다.
% It performs the following functions:
이 계층은 다음과 같은 기능을 수행합니다:
\begin{itemize}
% \item \textbf{Checklist Integration}: It ingests a master inspection checklist (Excel format) that defines the targets, expected states, and spatial metadata (e.g., \texttt{rear} tags, grid positions).
    \item \textbf{체크리스트 통합}: 대상, 예상 상태 및 공간 메타데이터(예: \texttt{rear} 태그, 그리드 위치)를 정의하는 마스터 점검 체크리스트(Excel 형식)를 수집합니다.
% \item \textbf{Task Dispatching}: Based on the equipment type defined in the checklist, it dispatches cropped image regions to the appropriate specialized inspector.
    \item \textbf{작업 디스패칭}: 체크리스트에 정의된 장비 유형에 따라 잘린 이미지 영역을 적절한 특화 점검기에 배정합니다.
% \item \textbf{Result Aggregation}: It collects diagnostic outputs from all modules, performs cross-validation (e.g., temporal consistency), and generates the final structured report.
    \item \textbf{결과 집계}: 모든 모듈에서 진단 출력을 수집하고, 교차 검증(예: 시간적 일관성)을 수행하며, 최종 구조화된 보고서를 생성합니다.
\end{itemize}

% \subsection{Perception and Pre-processing Layer: Bridging Real-world Noise}
\subsection{인지 및 전처리 계층: 실제 환경의 노이즈 극복}
% This layer transforms raw sensor data into refined inputs for diagnosis.
이 계층은 원시 센서 데이터를 진단을 위한 정제된 입력으로 변환합니다.
% It addresses the inherent noise and geometric distortions of industrial environments:
산업 환경 고유의 노이즈와 기하학적 왜곡을 처리합니다:
\begin{itemize}
% \item \textbf{Detection and Localization}: Employs an Ultralytics YOLOv11 model to detect primary equipment (gauges, switches) and secondary features (needles, scale marks).
    \item \textbf{탐지 및 지역화}: Ultralytics YOLOv11 모델을 사용하여 주요 장비(게이지, 스위치) 및 보조 특징(바늘, 눈금 표시)을 탐지합니다.
% \item \textbf{Geometric Refinement}: A crucial step where Homography-based projection is used to rectify perspective distortions, and Hough transform-based skew correction is applied to display screens.
    \item \textbf{기하학적 정제}: 호모그래피 기반 투영을 사용하여 원근 왜곡을 교정하고, 디스플레이 화면에 허프 변환 기반 스큐 교정을 적용하는 중요한 단계입니다.
% \item \textbf{Environmental Adaptation}: Includes adaptive CLAHE for low-light enhancement and specular reflection masking to stabilize visual features.
    \item \textbf{환경 적응}: 저조도 개선을 위한 적응형 CLAHE와 시각적 특징을 안정화하기 위한 정반사 마스킹을 포함합니다.
\end{itemize}

% \subsection{Specialized Diagnostics Layer: Domain-specific Interpreters}
\subsection{특화 진단 계층: 도메인 특화 해석기}
% Once the image regions are localized and rectified, they are processed by specialized engines:
이미지 영역이 지역화되고 교정되면 특화된 엔진에 의해 처리됩니다:
\begin{itemize}
% \item \textbf{Analog Gauge (AG) Inspector}: Uses pose estimation to identify the center point, scale bounds, and needle vector. The reading is calculated via geometric angle interpolation.
    \item \textbf{아날로그 게이지(AG) 점검기}: 포즈 추정을 사용하여 중심점, 눈금 범위 및 바늘 벡터를 식별합니다. 판독값은 기하학적 각도 보간을 통해 계산됩니다.
% \item \textbf{Digital Gauge (DG) Inspector}: Utilizes PaddleOCR to extract textual data, with post-processing logic to handle decimal points and multi-segment artifacts.
    \item \textbf{디지털 게이지(DG) 점검기}: PaddleOCR을 사용하여 텍스트 데이터를 추출하며, 소수점 및 멀티 세그먼트 아티팩트를 처리하기 위한 후처리 로직을 포함합니다.
% \item \textbf{Switch/LED Inspector}: Performs state classification (ON/OFF/COLOR) using a combination of color-space thresholding and delimiter-aware prefix matching against the checklist.
    \item \textbf{스위치/LED 점검기}: 색 공간 임계값 처리와 체크리스트에 대한 구분자 인식 접두사 매칭의 조합을 사용하여 상태 분류(ON/OFF/COLOR)를 수행합니다.
% \item \textbf{VLM Reasoning Engine}: Triggered for "ETC" items, this module utilizes a local Qwen2-VL model to perform semantic reasoning on non-standard safety equipment and complex nameplates.
    \item \textbf{VLM 추론 엔진}: "ETC" 항목에 대해 트리거되는 이 모듈은 로컬 Qwen2-VL 모델을 활용하여 비표준 안전 장비 및 복잡한 명판에 대한 의미론적 추론을 수행합니다.
\end{itemize}

% \subsection{Resource Management Layer: Performance and Scaling}
\subsection{자원 관리 계층: 성능 및 확장성}
% To ensure the system remains viable for real-time robotic operations, we implement several optimization strategies:
로봇의 실시간 운영이 가능하도록 시스템을 유지하기 위해 몇 가지 최적화 전략을 구현합니다:
\begin{itemize}
% \item \textbf{Model Management}: The \texttt{ModelManager} implements a singleton-pattern cache to prevent redundant model loading and manage VRAM allocation effectively.
    \item \textbf{모델 관리}: \texttt{ModelManager}는 중복 모델 로딩을 방지하고 VRAM 할당을 효과적으로 관리하기 위해 싱글톤 패턴 캐시를 구현합니다.
% \item \textbf{Parallel Execution}: A \texttt{ThreadPoolExecutor} is utilized to handle computationally intensive VLM requests in parallel, significantly reducing the bottleneck of semantic reasoning.
    \item \textbf{병렬 실행}: \texttt{ThreadPoolExecutor}를 활용하여 계산 집약적인 VLM 요청을 병렬로 처리함으로써 의미론적 추론의 병목 현상을 크게 줄입니다.
% \item \textbf{Temporal Consistency}: To eliminate transient errors, a voting mechanism operates across a sliding temporal window of three frames.
    \item \textbf{시간적 일관성}: 일시적인 오류를 제거하기 위해 3개 프레임의 슬라이딩 시간 윈도우에 걸쳐 투표 메커니즘이 작동합니다.
\end{itemize}
